\section{Exercise 2: AM Simulation}
\label{sec:lab2_ex2}

\subsection{Objective}

Integrate the AM modulator and both demodulation sub-VIs into a single top-level
simulation and compare coherent versus envelope detection as the modulation index
changes.

\subsection{Top-Level Design}

The top-level VI:
\begin{itemize}
\item generates an AM waveform from message and carrier parameters,
\item passes the same AM signal to coherent and envelope demodulators in
parallel,
\item displays time-domain and PSD outputs for direct comparison,
\item runs continuously so parameter changes can be observed in real time.
\end{itemize}

\begin{figure}[H]
\centering
\labincludegraphics[width=\linewidth]{figures/AM_top_level.png}
\caption{Top-level AM simulation VI connecting modulation, coherent demodulation, and envelope detection.}
\label{fig:lab2_am_top_level}
\end{figure}

\subsection{Parameter Set Used}

To match the handout configuration, the simulation uses:
\begin{itemize}
\item carrier amplitude $A_c = 2$,
\item message amplitude $A_m \in \{1,2,3,4\}$,
\item sample frequency $f_s = 200~\text{kHz}$,
\item number of samples $N = 1000$,
\item carrier frequency $f_c = 10~\text{kHz}$,
\item message frequency $f_m = 1~\text{kHz}$,
\item low-pass filter order $= 5$,
\item low-pass cutoff $= 3~\text{kHz}$.
\end{itemize}

The corresponding modulation indices are
\begin{equation}
\mu=\frac{A_m}{A_c}\in\{0.5,1.0,1.5,2.0\}.
\end{equation}

\begin{figure}[H]
\centering
\begin{subfigure}{0.49\linewidth}
\centering
\labincludegraphics[width=\linewidth]{figures/ex2_amplitude_1.png}
\caption{AM simulation result for $A_m = 1$ ($\mu = 0.5$).}
\label{fig:lab2_ex2_am1}
\end{subfigure}
\hfill
\begin{subfigure}{0.49\linewidth}
\centering
\labincludegraphics[width=\linewidth]{figures/ex2_amplitude_2.png}
\caption{AM simulation result for $A_m = 2$ ($\mu = 1.0$).}
\label{fig:lab2_ex2_am2}
\end{subfigure}

\vspace{0.5em}

\begin{subfigure}{0.49\linewidth}
\centering
\labincludegraphics[width=\linewidth]{figures/ex2_amplitude_3.png}
\caption{AM simulation result for $A_m = 3$ ($\mu = 1.5$).}
\label{fig:lab2_ex2_am3}
\end{subfigure}
\hfill
\begin{subfigure}{0.49\linewidth}
\centering
\labincludegraphics[width=\linewidth]{figures/ex2_amplitude_4.png}
\caption{AM simulation result for $A_m = 4$ ($\mu = 2.0$).}
\label{fig:lab2_ex2_am4}
\end{subfigure}
\caption{Top-level simulation captures used to compare coherent and envelope detection as message amplitude changes.}
\label{fig:lab2_ex2_all_cases}
\end{figure}

\subsection{What Changes Between the Two Demodulation Techniques? Why?}

As the message amplitude increases, the coherent demodulated sinusoid increases
in amplitude approximately in proportion to $A_m/2$, which is consistent with
the coherent detection result
\begin{equation}
y_{\mathrm{LPF}}(t)\approx \frac{1}{2}\left[A_c + A_m\cos(2\pi f_m t)\right]
\end{equation}
under synchronized carrier recovery.

The main difference between the two demodulators is the waveform shape:
\begin{itemize}
\item \textbf{coherent demodulation} preserves the signed baseband waveform after
DC removal,
\item \textbf{envelope detection} preserves the magnitude variation of the
envelope, so the recovered output tends to be positive-biased before DC removal.
\end{itemize}

\subsection{Effect of Modulation Index}

For $\mu = 0.5$ and $\mu = 1.0$, envelope detection works well because the
envelope does not invert. The coherent demodulator also works in both cases and
produces a recovered sinusoid with the expected amplitude scaling.

For $\mu = 1.5$ and $\mu = 2.0$, the system is over-modulated. Coherent
demodulation still recovers the message because it does not rely on envelope
shape alone, but the envelope detector becomes distorted because rectification
cannot recover the original message once the envelope has inverted.

\subsection{Discussion}

This simulation confirms the main theoretical difference between the two
demodulation techniques. Coherent detection is more faithful but requires
carrier synchronization. Envelope detection is simpler, but its usable operating
range is limited by the condition $\mu \le 1$.
